% ─────────────────────────────────────────────────────────────────────────────
\section{Introduction}
% ─────────────────────────────────────────────────────────────────────────────

Search and rescue operations are extremely time-critical because the probability of survival decreases rapidly as time passes. Ground teams are often constrained by difficult terrain and face the same hazards as the victims they seek. Unmanned Aerial Vehicles (UAVs) offer a superior alternative as they can be airborne in minutes and traverse large areas at high speeds regardless of the ground conditions.

The central challenge for any multi-UAV mission is effective coordination. Traditional solutions often rely on a central station to partition the search space and assign tasks. However, this architecture is fragile in disaster zones where communication is unreliable and the central controller represents a single point of failure. Decentralized multi-agent systems offer a more resilient approach by using local rules and shared environmental cues. This method is inspired by social insects and remains robust even if individual agents fail as it does not require a global view \citep{bk:Bonabeau1999}.

This paper presents a compact model of decentralized UAV search and rescue implemented in AnyLogic. Our model uses two main design principles. First, UAVs coordinate through a reversed Ant Colony Optimization (ACO) mechanism where pheromones mark covered ground to repel future agents. This guides the swarm toward unexplored territory. Second, agent logic is defined by a statechart inspired by the Belief-Desire-Intention (BDI) architecture \citep{bk:Wooldridge2009}. This makes behavioral modes and transition conditions explicit and easy to analyze.

Our contributions include a complete agent-based model with three interacting agent types. We formalize the scoring function for reversed ACO and explain how it governs the balance between exploration and exploitation. Finally, we establish an automated experimental pipeline to characterize how movement rules and behavioral parameters shape global search performance.
